% ═══════════════════════════════════════════════════════════════
% LH-02d: Lehrerhinweise Wiederholung + Stundenleistung Mi familia
% Spanisch Klasse 7 - Sequenz 2
% ═══════════════════════════════════════════════════════════════

\documentclass[11pt, a4paper]{article}

\newif\ifswmodus
\swmodusfalse

\usepackage[utf8]{inputenc}
\usepackage[T1]{fontenc}
\usepackage[ngerman]{babel}
\usepackage{helvet}
\renewcommand{\familydefault}{\sfdefault}

\usepackage[margin=2cm]{geometry}
\usepackage{enumitem}
\usepackage{tabularx}
\usepackage{booktabs}
\usepackage{longtable}

\usepackage{xcolor}
\usepackage{amssymb}
\usepackage{tcolorbox}
\tcbuselibrary{skins}

\usepackage{fancyhdr}
\usepackage{lastpage}
\usepackage{needspace}

% FARBEN
\definecolor{spanisch-color}{HTML}{c41e3a}
\definecolor{grau-color}{HTML}{888888}
\definecolor{hellgrau-color}{HTML}{f5f5f5}
\definecolor{basis-color}{HTML}{2a7a4b}
\definecolor{standard-color}{HTML}{2c5aa0}
\definecolor{erweiterung-color}{HTML}{7b1fa2}
\definecolor{loesung-color}{HTML}{c62828}

\definecolor{sw-dark}{HTML}{000000}
\definecolor{sw-gray}{HTML}{666666}
\definecolor{sw-white}{HTML}{ffffff}

\ifswmodus
    \colorlet{spanisch}{sw-dark}
    \colorlet{grau}{sw-gray}
    \colorlet{hellgrau}{sw-white}
    \colorlet{basis}{sw-dark}
    \colorlet{standard}{sw-dark}
    \colorlet{erweiterung}{sw-dark}
    \colorlet{loesung}{sw-dark}
\else
    \colorlet{spanisch}{spanisch-color}
    \colorlet{grau}{grau-color}
    \colorlet{hellgrau}{hellgrau-color}
    \colorlet{basis}{basis-color}
    \colorlet{standard}{standard-color}
    \colorlet{erweiterung}{erweiterung-color}
    \colorlet{loesung}{loesung-color}
\fi

\newcommand{\fachfarbe}{spanisch}

\newcommand{\thema}{Wiederholung + Stundenleistung Mi familia}
\newcommand{\sequenz}{02}
\newcommand{\block}{d}
\newcommand{\fach}{Spanisch}
\newcommand{\klasse}{7}
\newcommand{\dauer}{90 Min. (Doppelstunde)}
\newcommand{\lerngruppengroesse}{24}
\newcommand{\datum}{\today}

\pagestyle{fancy}
\fancyhf{}
\renewcommand{\headrulewidth}{0.4pt}
\renewcommand{\footrulewidth}{0.4pt}

\lhead{\fach{} -- Klasse \klasse}
\chead{\textbf{LH-\sequenz\block{} (Lehrerhinweise)}}
\rhead{\datum}

\lfoot{}
\cfoot{}
\rfoot{Seite \thepage\ von \pageref{LastPage}}

% Differenzierungsmarker
\newcommand{\B}{\textcolor{basis}{\textbf{[B]}}}
\newcommand{\Std}{\textcolor{standard}{\textbf{[S]}}}
\newcommand{\E}{\textcolor{erweiterung}{\textbf{[E]}}}

\newcommand{\loesung}[1]{\textcolor{loesung}{\textbf{#1}}}
\newcommand{\es}[1]{\textcolor{\fachfarbe}{\textbf{#1}}}

% Boxen
\newtcolorbox{kurzplan}{
    colback=hellgrau,
    colframe=\fachfarbe,
    colbacktitle=white,
    coltitle=\fachfarbe,
    boxrule=1pt,
    arc=0pt,
    fonttitle=\bfseries,
    attach boxed title to top left={yshift=-2mm, xshift=5mm},
    boxed title style={boxrule=0pt, colback=white},
    title=Kurzplan
}

\newtcolorbox{differenzierung}{
    colback=white,
    colframe=grau,
    colbacktitle=white,
    coltitle=grau,
    boxrule=0.5pt,
    arc=0pt,
    fonttitle=\bfseries,
    attach boxed title to top left={yshift=-2mm, xshift=5mm},
    boxed title style={boxrule=0pt, colback=white},
    title=Differenzierung
}

\newtcolorbox{fehlerbox}{
    colback=white,
    colframe=loesung,
    colbacktitle=white,
    coltitle=loesung,
    boxrule=0.5pt,
    arc=0pt,
    fonttitle=\bfseries,
    attach boxed title to top left={yshift=-2mm, xshift=5mm},
    boxed title style={boxrule=0pt, colback=white},
    title=Häufige Fehler
}

\newtcolorbox{materialbox}{
    colback=white,
    colframe=\fachfarbe,
    colbacktitle=white,
    coltitle=\fachfarbe,
    boxrule=0.5pt,
    arc=0pt,
    fonttitle=\bfseries,
    attach boxed title to top left={yshift=-2mm, xshift=5mm},
    boxed title style={boxrule=0pt, colback=white},
    title=Material-Checkliste
}

\newtcolorbox{testbox}{
    colback=white,
    colframe=\fachfarbe,
    colbacktitle=white,
    coltitle=\fachfarbe,
    boxrule=1pt,
    arc=0pt,
    fonttitle=\bfseries,
    attach boxed title to top left={yshift=-2mm, xshift=5mm},
    boxed title style={boxrule=0pt, colback=white},
    title=Stundenleistung SL-02d
}

\begin{document}

% ═══════════════════════════════════════════════════════════════
% KOPFBEREICH
% ═══════════════════════════════════════════════════════════════

\begin{center}
    {\Large\bfseries\textcolor{\fachfarbe}{Lehrerhinweise: \thema}}\\[0.3em]
    {\normalsize LH-\sequenz\block{} | \dauer{} | Sequenz 2 ABSCHLUSS}
\end{center}

\vspace{10pt}

% ═══════════════════════════════════════════════════════════════
% 1. KURZPLAN
% ═══════════════════════════════════════════════════════════════

\section*{1. Stundenverlauf (Kurzplan)}

\textbf{1. Stunde (45 Min.) -- Wiederholung + Übung}

\begin{tabularx}{\textwidth}{|c|l|X|l|}
\hline
\textbf{Zeit} & \textbf{Phase} & \textbf{Inhalt} & \textbf{Material} \\
\hline
0--5 & Einstieg & Begrüßung, Ankündigung: Wiederholung + Test & -- \\
\hline
5--12 & Wiederholung 1 & Familienwortschatz abfragen (mündlich) & Tafel \\
\hline
12--20 & Wiederholung 2 & tener-Konjugation, Possessivbegleiter & Tafel \\
\hline
20--30 & Übung AB & AB-02d Aufgaben 1-2 bearbeiten & AB \\
\hline
30--40 & Übung AB & AB-02d Aufgaben 3-4 bearbeiten & AB \\
\hline
40--45 & Sicherung & Lösungen besprechen, letzte Fragen & Tafel \\
\hline
\end{tabularx}

\vspace{1em}

\textbf{2. Stunde (45 Min.) -- Stundenleistung}

\begin{tabularx}{\textwidth}{|c|l|X|l|}
\hline
\textbf{Zeit} & \textbf{Phase} & \textbf{Inhalt} & \textbf{Material} \\
\hline
0--5 & Aktivierung & Letzte mündliche Wiederholung: Beispielsätze & -- \\
\hline
5--10 & Vorbereitung & SL-02d erklären, Fragen klären, Ruhe & SL \\
\hline
10--35 & Test & Stundenleistung (25 Min.), Aufsicht & SL \\
\hline
35--40 & Abgabe & Einsammeln, kurze Entspannung & -- \\
\hline
40--45 & Abschluss & Ausblick Sequenz 3 (Mi casa), Feedback & -- \\
\hline
\end{tabularx}

% ═══════════════════════════════════════════════════════════════
% 2. DIFFERENZIERUNG
% ═══════════════════════════════════════════════════════════════

\section*{2. Differenzierung}

\begin{differenzierung}
\textbf{Legende:} \B{} Basis (BR) | \Std{} Standard (MR) | \E{} Erweiterung (GY)

\vspace{0.5em}

\textbf{WICHTIG:} Diese Marker sind NUR für die Lehrkraft. Auf Schülermaterial erscheinen KEINE Niveaumarkierungen!
\end{differenzierung}

\vspace{1em}

\begin{tabularx}{\textwidth}{|l|X|X|X|}
\hline
& \B{} \textbf{Basis} & \Std{} \textbf{Standard} & \E{} \textbf{Erweiterung} \\
\hline
\textbf{AB Aufg. 1} & Mit Wortliste & Selbstständig & + Sätze bilden \\
\hline
\textbf{AB Aufg. 2} & Formen vorgegeben & Selbstständig & + Eigene Sätze \\
\hline
\textbf{AB Aufg. 3} & Mit Erklärung & Selbstständig & + Erklärung Regel \\
\hline
\textbf{AB Aufg. 4} & 3 Sätze genügen & 4+ Sätze & 6+ Sätze, komplex \\
\hline
\textbf{SL} & Aufg. 1-3 Pflicht (14P) & Alle Aufgaben (20P) & + kreative Extras \\
\hline
\end{tabularx}

\vspace{1em}

\textbf{Konkrete Hinweise:}
\begin{itemize}
    \item \B{} SuS mit LRS: Zeitzugabe (+5 Min.), Wortliste bei SL erlaubt
    \item \B{} DaZ-SuS: Familienbegriffe deutsch-spanisch vorab besprechen
    \item \E{} Leistungsstarke: Können anderen helfen, Zusatztext schreiben
\end{itemize}

% ═══════════════════════════════════════════════════════════════
% 3. HÄUFIGE FEHLER
% ═══════════════════════════════════════════════════════════════

\section*{3. Häufige Fehler}

\begin{fehlerbox}
\begin{tabularx}{\textwidth}{|l|X|X|}
\hline
\textbf{Fehler} & \textbf{Beispiel} & \textbf{Reaktion} \\
\hline
tener/ser verwechselt & *``Yo soy una hermana'' & Bedeutungsunterschied: haben vs. sein \\
\hline
Genus Possessiv falsch & *``mi padre'' statt ``mi madre'' & Possessiv richtet sich nach Objekt! \\
\hline
Plural vergessen & *``mi hermanos'' & mi $\rightarrow$ mis bei Plural \\
\hline
tengo/tiene verwechselt & *``Ella tengo'' & Personalpronomen betonen \\
\hline
Akzent vergessen & *``numero'' statt ``número'' & Bei Bewertung tolerieren wenn Bedeutung klar \\
\hline
\end{tabularx}
\end{fehlerbox}

% ═══════════════════════════════════════════════════════════════
% 4. MATERIAL-CHECKLISTE
% ═══════════════════════════════════════════════════════════════

\section*{4. Material-Checkliste}

\begin{materialbox}
\begin{itemize}[label=$\square$]
    \item AB-02d.pdf (\lerngruppengroesse{} Kopien)
    \item ML-02d.pdf (1× für Lehrkraft)
    \item SL-02d.pdf (\lerngruppengroesse{} Kopien -- STUNDENLEISTUNG)
    \item SL-02d-ML.pdf (1× für Lehrkraft)
    \item Lehrwerk: Qué pasa 1, S. 54-69
    \item Tafelmarker / Kreide
    \item Optional: LP-02d.html (Tablets/PCs für Wiederholung)
    \item Stoppuhr / Timer für Testzeit
\end{itemize}
\end{materialbox}

% ═══════════════════════════════════════════════════════════════
% 5. STUNDENLEISTUNG ÜBERSICHT
% ═══════════════════════════════════════════════════════════════

\section*{5. Stundenleistung SL-02d}

\begin{testbox}
\textbf{Dauer:} 25 Minuten | \textbf{Punkte:} 20 | \textbf{Hilfsmittel:} keine

\vspace{1em}

\begin{tabular}{|c|l|c|c|l|}
\hline
\textbf{Aufg.} & \textbf{Inhalt} & \textbf{AFB} & \textbf{Punkte} & \textbf{Hinweis} \\
\hline
1 & Zuordnung Familienmitglieder & I & 4 & 1P pro Zuordnung \\
\hline
2 & tener einsetzen & I/II & 4 & 1P pro Form \\
\hline
3 & Possessiv + Zahlen & II & 6 & 1P pro Lücke \\
\hline
4 & Dialog vervollständigen & II/III & 6 & 2P pro Antwort \\
\hline
\multicolumn{3}{|r|}{\textbf{Gesamt:}} & \textbf{20} & \\
\hline
\end{tabular}
\end{testbox}

\vspace{1em}

\textbf{Notenschlüssel (Klasse 7):}

\begin{tabular}{|c|c|c|c|c|c|c|}
\hline
\textbf{Note} & 1 & 2 & 3 & 4 & 5 & 6 \\
\hline
\textbf{ab \%} & 96 & 80 & 60 & 40 & 20 & <20 \\
\hline
\textbf{ab Punkte} & 19 & 16 & 12 & 8 & 4 & <4 \\
\hline
\end{tabular}

% ═══════════════════════════════════════════════════════════════
% 6. LÖSUNGEN AB-02d
% ═══════════════════════════════════════════════════════════════

\section*{6. Lösungen AB-02d}

\textbf{Merkkasten:}
\begin{tabular}{ll}
la madre & = die Mutter \\
el padre & = der Vater \\
la hermana & = die Schwester \\
el hermano & = der Bruder \\
la abuela & = die Großmutter \\
el abuelo & = der Großvater \\
\end{tabular}

\vspace{1em}

\textbf{Aufgabe 1 (Zuordnung):} (4P)
\begin{tabular}{ll}
\es{la madre} & $\rightarrow$ \loesung{die Mutter} \\
\es{el padre} & $\rightarrow$ \loesung{der Vater} \\
\es{la hermana} & $\rightarrow$ \loesung{die Schwester} \\
\es{el abuelo} & $\rightarrow$ \loesung{der Großvater} \\
\end{tabular}

\vspace{1em}

\textbf{Aufgabe 2 (tener):} (4P)
\begin{tabular}{ll}
a) Yo & \loesung{tengo} \\
b) Tú & \loesung{tienes} \\
c) Ella & \loesung{tiene} \\
d) Nosotros & \loesung{tenemos} \\
\end{tabular}

\vspace{1em}

\textbf{Aufgabe 3 (Possessivbegleiter):} (6P)
\begin{tabular}{ll}
a) & \loesung{Mi} madre \\
b) & \loesung{tu} hermano \\
c) & \loesung{Su} padre \\
d) & \loesung{Nuestra} abuela \\
e) & \loesung{Mis} hermanos \\
f) & \loesung{Sus} primos \\
\end{tabular}

\vspace{1em}

\textbf{Aufgabe 4 (Text):} (6P) -- Beispiellösung siehe ML-02d

% ═══════════════════════════════════════════════════════════════
% 7. LÖSUNGEN SL-02d
% ═══════════════════════════════════════════════════════════════

\section*{7. Lösungen SL-02d (Stundenleistung)}

Siehe separate Datei: \textbf{SL-02d-ML.pdf}

% ═══════════════════════════════════════════════════════════════
% 8. HAUSAUFGABE
% ═══════════════════════════════════════════════════════════════

\section*{8. Hausaufgabe}

\begin{itemize}
    \item Vokabeln Sequenz 2 wiederholen für Vokabeltest
    \item Optional: LP-02d.html zur Wiederholung (vor dem Test)
\end{itemize}

% ═══════════════════════════════════════════════════════════════
% 9. REFLEXIONSNOTIZEN
% ═══════════════════════════════════════════════════════════════

\section*{9. Reflexionsnotizen (nach Durchführung)}

\begin{tabularx}{\textwidth}{|l|X|}
\hline
\textbf{Was lief gut?} & \\[2em]
\hline
\textbf{Was lief nicht?} & \\[2em]
\hline
\textbf{Zeitmanagement} & \\[2em]
\hline
\textbf{Test-Ergebnisse} & Durchschnitt: \hspace{2cm} Beste Note: \hspace{2cm} Schlechteste: \\[1em]
\hline
\textbf{Für nächstes Mal} & \\[2em]
\hline
\end{tabularx}

\end{document}
